\documentclass[12pt]{article}

\usepackage[a4paper, margin=1in]{geometry}

\usepackage[czech]{babel}
\usepackage[utf8]{inputenc}
\usepackage[T1]{fontenc}
\usepackage{lingmacros}
\usepackage{tree-dvips}
\usepackage{listings}
\usepackage{xstring}
\usepackage{xparse}
\usepackage{booktabs}
\usepackage{longtable}
\usepackage[hidelinks]{hyperref}
\usepackage{parskip}
\usepackage{titling}

\author{\textbf{Tomáš Tobrman}}
\title{Semestrální práce \textbf{KIV/ZOS}}
\date{}

\renewcommand*\contentsname{Obsah}
\renewcommand\maketitlehooka{\null\mbox{}\vfill}
\renewcommand\maketitlehookd{\vfill\null}

\begin{document}

\begin{titlingpage}
\maketitle
\thispagestyle{empty}
\end{titlingpage}
\newpage

\tableofcontents
\thispagestyle{empty}
\newpage

\setcounter{page}{1}

\section{Úvod}

\subsection{Zadání}

Vytvoření konzolového rozhraní s~příkazy pro práci s~pseudosystémem založeným na souborovém systému FAT.

\subsection{Implementace}

Projekt byl implementován v~programovacím jazyce C~pro operační systémy Linux/Unix.
Každému příkazu připadá vlastní zdrojový soubor, ve kterém se nachází jeho funkcionalita.
Pole příkazů je pak definováno ve zdrojovém souboru \texttt{utils.c}, proti kterému se porovnává uživatelský vstup.

\section{Použití}\label{sec:usage}

Před spuštěním programu je třeba zahájit jeho kompilaci.
Na systémech Linux/Unix je možné tohoto docílit pomocí příkazu \texttt{make}.
Po úspěšném sestavení program lze spustit z~příkazové řádky ve složce \texttt{bin}.

\subsection{Příkazy}\label{sec:commands}

Tato sekce obsahuje popis a~příklady využití všech příkazů implementovaných v~konzolovém rozhraní.

\subsubsection{\texttt{cat}}\label{sec:cat}

Vypíše obsah souboru.

\texttt{\$ cat file}

Příkaz očekává vždy jeden argument, který je cestou k~souboru, jehož obsah chce uživatel vypsat.

\subsubsection{\texttt{cd}}\label{sec:cd}

Změní cestu aktuálního pracovního adresáře.

\texttt{\$ cd path}\qquad
\texttt{\$ cd /}\qquad
\texttt{\$ cd}

Příkaz očekává až jeden argument, který odpovídá cestě do adresáře.
Pokud není pos\-kytnuta cesta, pracovní adresář je změněn na domácí složku uživatele (v~tomto případě kořenový adresář \texttt{/}).

\subsubsection{\texttt{clear}}\label{sec:clear}

Vyčistí obrazovku terminálu a~položí kurzor na první řádek.

\texttt{\$ clear}

Příkaz neočekává žádné argumenty.

\subsubsection{\texttt{cp}}\label{sec:cp}

Vytvoří kopii souboru.

\texttt{\$ cp source destination}

Příkaz očekává právě dva argumenty.
První specifikuje cestu/název souboru, který má být zkopírovaný.
Druhý argument udává cestu destinace a~název kopie.

\subsubsection{\texttt{defrag}}\label{sec:defrag}

Provádí defragmentaci specifikovaného souboru.

\texttt{\$ defrag file}

Příkaz očekává právě jeden argument, kterým je cesta/název souboru určeného pro defragmentaci.

\subsubsection{\texttt{exit}}\label{sec:exit}

Ukončí program.

\texttt{\$ exit}\qquad
\texttt{\$ exit code}

Při spuštění příkazu \texttt{exit} bez argumentu, program je ukončen s~návratovou hodnotou \texttt{0}.
Použitím volitelného argumentu lze specifikovat jinou návratovou hodnotu programu (modulo 256).

\subsubsection{\texttt{format}}\label{sec:format}

Provádí formátování souborového systému.
Tento příkaz bývá prvním, který musí být proveden.

\texttt{\$ format size}

Příkaz očekává právě jeden argument, který udává velikost souborového systému.

Očekávané jednotky jsou:

\texttt{B} (Byte),
\texttt{kB} (Kilobyte),
\texttt{MB} (Megabyte),
\texttt{GB} (Gygabyte)

Maximální velikost souborového systému je \texttt{1GB} a~minimální jsou \texttt{2kB}.

\subsubsection{\texttt{help}}\label{sec:help}

Vypíše základní informace o~všech příkazech.

\texttt{\$ help}

Příkaz neočekává žádné argumenty.

\newpage

\subsubsection{\texttt{incp}}\label{sec:incp}

Vytvoří kopii souboru z~hostitelského zařízení do virtuálního souborového systému.

\texttt{\$ incp source destination}

Funkce příkazu je stejná jako příkaz \texttt{cp}.
Zdrojový soubor se musí nacházet na hostitelském zařízení.

\subsubsection{\texttt{info}}\label{sec:info}

Vypíše informace o~adresáři nebo souboru.

\texttt{\$ info item}

Příkaz očekává právě jeden argument, kterým je cesta/název souboru či adresáže.

\subsubsection{\texttt{load}}\label{sec:load}

Načte soubor z~hostitelského zařízení a~provede v~něm vypsané příkazy.

\texttt{\$ load file}

Příkaz očekává právě jeden argument specifikující cestu/název souboru na hostitelském zařízení.

\subsubsection{\texttt{ls}}\label{sec:ls}

Vypíše obsah adresáře.

\texttt{\$ ls}\qquad
\texttt{\$ ls path}\qquad
\texttt{\$ ls -l path}\qquad
\texttt{\$ ls -a path}\qquad
\texttt{\$ ls -la path}

Pokud je příkaz proveden bez argumentů, vypíše viditelné položky aktuálního adresáře.
Je možné poskytnout argument cesty adresáře pro vypsání.

Přepínač \texttt{-l} vypíše obsah adresáře v~seznamovém formátu.

Přepínač \texttt{-a} vypíše také schované položky (začínající tečkou).

\subsubsection{\texttt{mkdir}}\label{sec:mkdir}

Vytvoří nový adresář.

\texttt{\$ mkdir path}

Příkaz očekává právě jeden argument s~cestou/názvem nového adresáře.

\subsubsection{\texttt{mv}}\label{sec:mv}

Přesune či přejmenuje soubor.

\texttt{\$ mv source destination}

Příkaz očekává právě dva argumenty. Tyto argumenty jsou stejné jako u~příkazu \texttt{cp}.

\subsubsection{\texttt{outcp}}\label{sec:outcp}

Vytvoří kopii souboru z~virtuálního souborového systému do hostitelského zařízení.

\texttt{\$ outcp source destination}

Funkce příkazu je stejná jako příkaz \texttt{cp}.
Cílový soubor bude vytvořen na hostitelském zařízení.

\subsubsection{\texttt{pwd}}\label{sec:pwd}

Vypíše cestu do aktuálního adresáře.

\texttt{\$ pwd}

Příkaz neočekává žádné argumenty.

\subsubsection{\texttt{rm}}\label{sec:rm}

Odstraní soubor ze souborového systému.

\texttt{\$ rm file}

Příkaz očekává právě jeden argument s~cestou/názvem souboru k~odstranění.

\subsubsection{\texttt{rmdir}}\label{sec:rmdir}

Odstraní prázdný adresář ze souborového systému.

\texttt{\$ rmdir directory}

Příkaz očekává právě jeden argument s~cestou/názvem prázdného adresáře k odstranění.

\section{Závěr}

Projekt byl vyvíjen a testován na operačním systému Linux.
Na jiné platformě se tedy mohou objevit problémy, které nebyly zastiženy při vývoji.

Uživatelské rozhraní programu je modelováno podle programu \texttt{bash}.
Příkaz \texttt{ls} navíc zobrazuje složky v modré barvě.
Pokud tedy Vaše konzole nepodporuje některé z funkcí, se kterými program počítá, je možné, že výstup nebude vypadat korektně.
Na funkci a běhu programu by však tyto chyby němely mít vliv.

\end{document}
